\documentclass[a4paper,12pt]{ltjarticle}
\usepackage[top=25truemm,bottom=25truemm,left=10truemm,right=10truemm]{geometry}
\usepackage{amsmath}
\usepackage{multirow}

\begin{document}

\section*{分散分析ワークシート}

このワークシートでは，実験データから効果の大きさや平方和を計算する練習をします。表の空欄を埋めながら，データの散らばりがどのように分解されるかを確認しましょう。

\subsection*{課題1：一要因3水準の分散分析}

ある実験で，3つの条件（A，B，C）に被験者を割り当て，課題の成績を測定しました。以下の表を完成させなさい。

\begin{table}[ht]
	\centering
	\caption{一要因3水準のデータ}
	\begin{tabular}{|c|c|c|c|c|c|c|}
		\hline
		被験者ID & 条件 & スコア & 全体平均 & 群平均 & 効果 & 誤差 \\ \hline
		1        & A  & 12   &      &      &    &    \\ \hline
		2        & A  & 14   &      &      &    &    \\ \hline
		3        & A  & 16   &      &      &    &    \\ \hline
		4        & A  & 18   &      &      &    &    \\ \hline
		5        & B  & 10   &      &      &    &    \\ \hline
		6        & B  & 12   &      &      &    &    \\ \hline
		7        & B  & 14   &      &      &    &    \\ \hline
		8        & B  & 16   &      &      &    &    \\ \hline
		9        & C  & 16   &      &      &    &    \\ \hline
		10       & C  & 18   &      &      &    &    \\ \hline
		11       & C  & 20   &      &      &    &    \\ \hline
		12       & C  & 22   &      &      &    &    \\ \hline
	\end{tabular}
\end{table}

\textbf{計算欄：}
\begin{itemize}
	\item[] 条件Aの平均：\underline{\hspace{3cm}}
	\item[] 条件Bの平均：\underline{\hspace{3cm}}
	\item[] 条件Cの平均：\underline{\hspace{3cm}}
	\item[] 全体平均：\underline{\hspace{3cm}}
\end{itemize}

\textbf{平方和の計算：}
\begin{itemize}
	\item[] 条件Aの効果：\underline{\hspace{3cm}}
	\item[] 条件Bの効果：\underline{\hspace{3cm}}
	\item[] 条件Cの効果：\underline{\hspace{3cm}}
	\item[] 条件間平方和（SS\_between）：\underline{\hspace{3cm}}
	\item[] 誤差平方和（SS\_error）：\underline{\hspace{3cm}}
	\item[] 全体平方和（SS\_total）：\underline{\hspace{3cm}}
\end{itemize}

\textbf{確認：}SS\_total = SS\_between + SS\_error が成り立つか確認しましょう。

\newpage

\subsection*{課題2：二要因（2×2）の分散分析}

ある実験で，2つの要因（要因A：2水準，要因B：2水準）を組み合わせた実験を行いました。以下の表を完成させなさい。

\begin{table}[ht]
	\centering
	\caption{二要因（2×2）のデータ}
	\small
	\begin{tabular}{|c|c|c|c|c|c|c|c|c|c|}
		\hline
		ID & 要因A & 要因B & スコア & 全体平均 & Aの効果 & Bの効果 & 交互作用 & 群平均 & 誤差 \\ \hline
		1        & $a_1$ & $b_1$ & 18   &        &          &          &        &      &    \\ \hline
		2        & $a_1$ & $b_1$ & 20   &        &          &          &        &      &    \\ \hline
		3        & $a_1$ & $b_1$ & 22   &        &          &          &        &      &    \\ \hline
		4        & $a_1$ & $b_2$ & 12   &        &          &          &        &      &    \\ \hline
		5        & $a_1$ & $b_2$ & 14   &        &          &          &        &      &    \\ \hline
		6        & $a_1$ & $b_2$ & 16   &        &          &          &        &      &    \\ \hline
		7        & $a_2$ & $b_1$ & 10   &        &          &          &        &      &    \\ \hline
		8        & $a_2$ & $b_1$ & 12   &        &          &          &        &      &    \\ \hline
		9        & $a_2$ & $b_1$ & 14   &        &          &          &        &      &    \\ \hline
		10       & $a_2$ & $b_2$ & 16   &        &          &          &        &      &    \\ \hline
		11       & $a_2$ & $b_2$ & 18   &        &          &          &        &      &    \\ \hline
		12       & $a_2$ & $b_2$ & 20   &        &          &          &        &      &    \\ \hline
	\end{tabular}
\end{table}

\textbf{計算欄：}
\begin{itemize}
	\item[] $a_1b_1$群の平均：\underline{\hspace{3cm}}
	\item[] $a_1b_2$群の平均：\underline{\hspace{3cm}}
	\item[] $a_2b_1$群の平均：\underline{\hspace{3cm}}
	\item[] $a_2b_2$群の平均：\underline{\hspace{3cm}}
	\item[] 全体平均：\underline{\hspace{3cm}}
\end{itemize}

\textbf{主効果の計算：}
\begin{itemize}
	\item[] $a_1$水準の平均：\underline{\hspace{3cm}} $\rightarrow$ 要因Aの効果（$a_1$）：\underline{\hspace{3cm}}
	\item[] $a_2$水準の平均：\underline{\hspace{3cm}} $\rightarrow$ 要因Aの効果（$a_2$）：\underline{\hspace{3cm}}
	\item[] $b_1$水準の平均：\underline{\hspace{3cm}} $\rightarrow$ 要因Bの効果（$b_1$）：\underline{\hspace{3cm}}
	\item[] $b_2$水準の平均：\underline{\hspace{3cm}} $\rightarrow$ 要因Bの効果（$b_2$）：\underline{\hspace{3cm}}
\end{itemize}

\textbf{交互作用の計算：}
\begin{itemize}
	\item[] $a_1b_1$の交互作用：\underline{\hspace{3cm}}
	\item[] $a_1b_2$の交互作用：\underline{\hspace{3cm}}
	\item[] $a_2b_1$の交互作用：\underline{\hspace{3cm}}
	\item[] $a_2b_2$の交互作用：\underline{\hspace{3cm}}
\end{itemize}

\textbf{平方和の計算：}
\begin{itemize}
	\item[] 要因Aの平方和（SS\_A）：\underline{\hspace{3cm}}
	\item[] 要因Bの平方和（SS\_B）：\underline{\hspace{3cm}}
	\item[] 交互作用の平方和（SS\_AB）：\underline{\hspace{3cm}}
	\item[] 誤差平方和（SS\_error）：\underline{\hspace{3cm}}
	\item[] 全体平方和（SS\_total）：\underline{\hspace{3cm}}
\end{itemize}

\textbf{確認：}SS\_total = SS\_A + SS\_B + SS\_AB + SS\_error が成り立つか確認しましょう。

\end{document}
